\section{Introduction}
\label{sec:introduction}

The appeal of activation sparsity is immediate: if an operand is exactly zero,
some multiply--accumulate work is mathematically unnecessary.  The usual
summary, however, is a zero percentage at a named tensor.  That percentage is
not a model-level computational quantity.  A zero before a wide MLP projection
has a different footprint from a zero in a query vector; a value zero can be
reused over many causal queries; a pre-RoPE query zero need not remain zero
after rotation; and a dense language-model head can dilute block sparsity.
Consequently, two methods with similar local sparsity can expose very different
amounts and locations of avoidable scalar products.

Prior work already recognizes this mismatch.  TEAL aggregates per-matrix
sparsity with matrix-footprint weights and applies training-free thresholding
through linear layers \citep{liu2024teal}.  Our question is narrower and
complementary: \emph{given the actual operands of a declared full-sequence
causal transformer graph, how many scalar products contain at least one exact
zero?}  We answer with $\Rmodel$, an integer-counted metric spanning fused QKV,
causal QK, causal PV, attention output, and both MLP projections, plus a dense
LM-head denominator.  This refinement matters wherever zeros are transformed
or reused between named activations and multiplications.  It is deliberately
agnostic to kernel support and hardware speed.

Measurement also changes the experimental question.  Rather than optimizing a
single site's sparsity, we compare validation loss against model-wide logical
opportunity.  We study post-hoc clipping and trained gates in Pythia models
constructed from architecture configurations with random initialization
\citep{biderman2023pythia}.  The trained recipes place one-sided gates at the
parallel attention and MLP branch inputs and MLP hidden state, optionally adding
post-RoPE query, key, value, and pre-output context gates.  An Adam-aware
orthogonal $\ell_1$ update encourages small activations while removing the
component that opposes the preconditioned task direction.

The results are not a simple ``more sparsity is better'' story.  First, the
operation-weighted composition is explanatory: the seven-site high-dose recipe
adds 10--17 percentage points from QK and PV alone, and can exceed the
four-site recipe despite lower zero mass at some local sites.  Second, scale
changes the observed trained response under our common-token protocol.  At 14M
and 70M, raising the gate threshold from 0 to 0.5 increases loss; at 410M it
reduces loss by 0.50 and 0.31 for the four- and seven-site recipes.  The
evaluation-only curves do not reverse.  Because all scales see the same 1.493B
tokens, the 410M model receives far less exposure per parameter and has a weak
baseline.  A focused higher-learning-rate screen does not fix it, but longer
training remains untested.  We retain this cohort because the disagreement
between trained and post-hoc behavior exposes an important boundary of the
current evidence.

Our contributions are:
\begin{itemize}[leftmargin=*,nosep]
  \item an exact, pooled-integer definition of model-wide zero-product
  opportunity for a declared causal transformer workload, including operand
  placement and reuse that per-matrix summaries can miss;
  \item a joint accounting of quality, sitewise zero structure, and
  per-operation contributions for post-hoc and trained sparsification recipes;
  \item a matched-token Pythia case study through 410M that identifies a
  reproducible small-model pattern and a scientifically useful boundary result,
  while explicitly separating logical opportunity from runtime and scaling
  claims.
\end{itemize}

